\begin{intuition}[Motivation]
  As seen in Corollary~\ref{thm:no-rational-square-root-two}, $\sqrt{2}$ is not
  in $\Q$: $\Q$ has ``holes.''

  Goal: define $\R$ as an ordered field $+$
  ``completeness axiom.''
\end{intuition}

\subsection{Maxima, Minima, and Bounds}

\begin{definition}{Maximum and Minimum}{maximum-minimum}
  Let $(\mathbb{F},+,\cdot,\le)$ be an ordered field,\\
  let $S\subseteq\mathbb{F}$ be a nonempty subset.
  \begin{enumerate}[label=(\roman*)]
    \item Say an element $s_0\in S$ is the \textbf{maximum},\\
          denoted $s_0=\max S$, if $\forall s\in S,\ s\le s_0$.
    \item Vice versa, $s_0\in S$ is the \textbf{minimum},\\
          $s_0=\min S$ if $\forall s\in S,\ s\ge s_0$.
  \end{enumerate}
\end{definition}

\begin{theorem}{Extrema of Finite Sets}{finite-set-extrema}
  $\forall$ finite $S\subseteq\mathbb{F}$, $S\ne\varnothing$, $S$ has
  max \& min.
\end{theorem}

\begin{theorem}{Extrema of the Natural Numbers}{natural-number-extrema}
  $\N\subseteq\Q$ has $\min=1$
  but no max.
\end{theorem}

\begin{theorem}{No Extrema in an Open Rational Interval}{open-rational-interval-extrema}
  $S=\{r\in\Q\mid0<r<1\}$ has no max or min.
\end{theorem}

\begin{proof}[Proof for the maximum]
  By contradiction, e.g. for max:\\
  assume $r_0=\max S$.
  \begin{gather*}
    r_1=\frac{r_0+1}{2}\in S\\
    \text{and }r_1>r_0.
  \end{gather*}
  $r_0$ does not exist.
\end{proof}

\begin{intuition}
  For $S$, $0,1$ feel like some edge?
\end{intuition}

\begin{definition}{Supremum and Infimum}{supremum-infimum}
  \textbf{Important!}\\
  For $(\mathbb{F},+,\cdot,\le)$ an ordered field,\\
  let $S\subseteq\mathbb{F}$, $S\ne\varnothing$.
  \begin{enumerate}[label=(\roman*)]
    \item An element $r\in\mathbb{F}$ is the \textbf{supremum},\\
          $r=\sup S$, if:
          \begin{enumerate}[label=(\alph*)]
            \item $r$ is an upper bound for $S$:\\
                  $s\le r,\ \forall s\in S$.
            \item $\forall t\in\mathbb{F}$ an upper bound for $S$,\\
                  $r\le t$.
          \end{enumerate}
          (\textbf{Least upper bound}.)
    \item Vice versa, $r\in\mathbb{F}$ is the \textbf{infimum} of $S$,\\
          $r=\inf S$: \textbf{greatest lower bound}.
  \end{enumerate}
\end{definition}

\begin{note}
  $\sup S$ and $\inf S$ are not necessarily in $S$.
\end{note}

\begin{example}
  Take $\mathbb{F}=\Q$. $S=\{r\in\Q\mid0<r<1\}$.
  \[
    \sup S=1,\qquad\inf S=0.
  \]
  \begin{proof}[Proof that $\sup S=1$]
    \begin{enumerate}[label=(\roman*)]
      \item $1$ is an upper bound by the definition of $S$.
      \item Let $t\in\Q$ be any other upper bound.\\
            Want to show $1\le t$.\\
            Prove the contrapositive: $1>t\implies t$ is not an upper bound.\\
            Borrow the argument from no max in $S$
            (Theorem~\ref{thm:open-rational-interval-extrema}).\\
            Such an upper bound $t<1$ must not exist.
    \end{enumerate}
    So $1$ is indeed the least upper bound.
  \end{proof}
\end{example}

\begin{theorem}{Extrema and Bounds}{extrema-and-bounds}
  For $S$,
  \begin{align*}
    \text{if }\exists\max S & \implies\max S=\sup S, \\
    \text{if }\exists\min S & \implies\min S=\inf S.
  \end{align*}
\end{theorem}

\textit{Proof:} $\ldots$

\begin{samepage}
  \begin{theorem}{Infimum Is at Most Supremum}{infimum-at-most-supremum}
    For $S\ne\varnothing$ with $\inf S,\sup S$ defined, $\sup S\ge\inf S$.
  \end{theorem}

  \begin{proof}
    \[
      S\ne\varnothing\implies\exists t\in S.
    \]
    $\inf S\le t$, $\sup S\ge t$ by Definition~\ref{def:supremum-infimum}.
    \[
      \inf S\le\sup S.\qedhere
    \]
  \end{proof}
\end{samepage}

\begin{theorem}{Equal Infimum and Supremum}{equal-infimum-supremum}
  For nonempty $S$ with an infimum and supremum,
  \[
    \inf S=\sup S\iff S=\{a\},\quad a\in\mathbb{F}.
  \]
\end{theorem}

\begin{proof}
  $(\Leftarrow)$ Trivial.

  $(\Rightarrow)$ Proof by contradiction:\\
  suppose $S$ has more than one element.
  \[
    \exists t_1,t_2\in S,\quad t_1\ne t_2.
  \]
  So $t_1<t_2$ or $t_1>t_2$.\\
  Let $t_1<t_2$.
  \[
    \inf S\le t_1<t_2\le\sup S,
  \]
  Not equal, contradiction.
\end{proof}

\subsection{Completeness of the Real Numbers}

\begin{example}
  For the set $S=\{r\in\Q\mid r^2<2\}$:\\
  bounded above, $\sup S$ is not in $\Q$!
\end{example}

\begin{theorem}{Existence of the Real Numbers}{existence-real-numbers}
  $\exists$ an ordered field $\R$, the \textbf{real numbers},
  satisfying the ``\textbf{completeness axiom}'':
  \[
    \begin{gathered}
      \text{If }S\subseteq\R,\ S\ne\varnothing\text{ and bounded above},\\
      \implies\exists\sup S\in\R.
    \end{gathered}
  \]
\end{theorem}

Hard to prove!

We view $\N\subseteq\Z\subseteq\Q\subseteq\R$.

\begin{intuition}
  Uniqueness of $\R$? Always a good thing to ask after existence.
\end{intuition}

\begin{theorem}{Uniqueness of the Real Numbers}{uniqueness-real-numbers}
  Formally, ordered fields $\mathbb{F}$ and $\R$ both satisfy the completeness axiom\\
  $\implies\mathbb{F}$ and $\R$ are isomorphic.
\end{theorem}

Hard to prove!

\begin{theorem}{Existence of Infima}{existence-infima}
  Let $S\subseteq\R$, $S\ne\varnothing$ and bounded below,\\
  then $\exists\inf S$.
\end{theorem}

\begin{proof}[Proof sketch]
  Consider $-S:=\{-s\mid s\in S\}$.\\
  Show $\sup(-S)$ exists\\
  and show $-\sup(-S)=\inf S$.
\end{proof}

\subsection{The Archimedean Property}

\begin{theorem}{Archimedean Property}{archimedean-property}
  Let $a,b\in\R$, s.t.\ $a,b>0$.
  \[
    \implies\exists n\in\N\text{ s.t. }na>b.
  \]
\end{theorem}

\begin{intuition}
  Weird things don't happen far out.
  \begin{center}
    \begin{tikzpicture}
      \draw[number line] (-0.5,0) -- (10,0);
      \foreach \x/\label in {0/0,1.2/a,2.4/2a,3.6/3a,7.7/b,9/na} {
        \draw (\x,-0.1) -- (\x,0.1) node[below=5pt] {$\label$};
      }
      \node[fill=orange!5,inner sep=4pt] at (5.7,0) {$\cdots$};
    \end{tikzpicture}
  \end{center}
\end{intuition}

\begin{note}
  \begin{enumerate}
    \item Proof for $\Q$ without completeness axiom.
    \item $\exists$ an ordered field where this fails
          (e.g. rational functions).
    \item It's a bit hard on $\R$ since we only know
          completeness axiom.
  \end{enumerate}
\end{note}

\begin{proof}
  On $\R$: by contradiction, assume
  \[
    \exists a,b>0\text{ s.t. }na\le b\quad\forall n\in\N.
  \]
  Let $S=\{na\mid n\in\N\}$.
  \begin{flushleft}
    $\implies b$ is an upper bound for $S$.\\
    $\implies$ by completeness (Theorem~\ref{thm:existence-real-numbers}),
    $\exists$ least upper bound $\exists\sup S$.
  \end{flushleft}
  Since $a>0$, $\sup S-a<\sup S$.
  \begin{align*}
     & \implies\sup S-a\text{ is not an upper bound}, \\
     & \implies\exists s\in S\text{ s.t. }\sup S-a<s.
  \end{align*}
  Since $s\in S$, we know $s=na$ for some $n\in\N$.
  \begin{align*}
     & \sup S-a<na        \\
     & \iff\sup S<(n+1)a.
  \end{align*}
  But $(n+1)a\in S$, $\sup S$ is not an upper bound.\\
  Contradiction!
\end{proof}

\begin{intuition}
  What if $b$ is so far away?\\[0.5em]
  If that's the case, $S=\{na\mid n\in\N\}$ is bounded,\\
  there is some limit $\sup S$, that does not make sense.
  \begin{center}
    \begin{tikzpicture}
      \draw[number line] (-0.5,0) -- (10.5,0);
      \foreach \x in {0.2,1.1,2} {\draw (\x,-0.08) -- (\x,0.08);}
      \node at (2.9,0.2) {$\cdots$};
      \draw (3.6,-0.1) -- (3.6,0.1) node[below=5pt] {$\sup S-a$};
      \node[sketch point,label=above:{$na$}] at (4.8,0) {};
      \draw[dashed] (6,-0.65) -- (6,0.85);
      \node[below] at (6,-0.65) {$\sup S$};
      \node[sketch point,label=below:{$(n+1)a$}] at (7.2,0) {};
      \draw (9.5,-0.1) -- (9.5,0.1) node[below=5pt] {$b$};
      \draw[->,>=latex,sketch accent] (4.8,0.25) to[bend left=40]
        node[above] {can go over} (7.2,0.25);
      \draw[decorate,decoration={brace,mirror,amplitude=4pt}]
        (3.6,-1.25) -- (6,-1.25) node[midway,below=6pt] {gap $a$};
      \node[below] at (6,-1.9) {assume bounded};
    \end{tikzpicture}
  \end{center}
\end{intuition}

\subsection{Density of the Rational Numbers}

\begin{theorem}{Density of $\Q$ in $\R$}{density-rationals}
  For $a,b\in\R$, $a<b$,
  \[
    \implies\exists r\in\Q,\text{ s.t. }a<r<b.
  \]
\end{theorem}

\begin{intuition}
  Weird things don't happen in $\R$ zoomed in.
\end{intuition}

\begin{flushleft}
  \textit{Proof sketch:} Want to show $a<\dfrac mn<b$.\\
  WTS $na<m<nb$.
\end{flushleft}

\begin{intuition}
  Take $n$ really big, that
  some integer finally lands here in the gap.
  \begin{center}
    \begin{tikzpicture}
      \draw (-0.2,0) -- (6.8,0);
      \foreach \x in {0,1.5,3} {
        \draw[->,>=latex,sketch accent] (\x,0.1) to[bend left=45] (\x+1.5,0.1);
      }
      \draw (3.4,-0.12) -- (3.4,0.5);
      \draw (6,-0.12) -- (6,0.65);
      \node[below=5pt] at (3.4,0) {$na$};
      \node[below=5pt] at (6,0) {$nb$};
      \node[below=5pt] at (4.5,0) {$m$};
      \draw[->,>=latex] (4.5,1) -- (4.5,0.15);
    \end{tikzpicture}
  \end{center}
\end{intuition}

\begin{proof}
  By the Archimedean property (Theorem~\ref{thm:archimedean-property}):
  \begin{gather*}
    \exists n\in\N,\text{ s.t. }n(b-a)>1,\\
    \iff nb-na>1.
  \end{gather*}
  Also by AP:
  \begin{gather*}
    \exists k\in\N,\text{ s.t. }k>\max\{\abs{na},\abs{nb}\},\\
    \implies-k<na<nb<k.
  \end{gather*}
  Let $m=\min\{j\in\Z\mid-k\le j\le k,\ na<j\}$.
  \begin{gather*}
    m-1\le na\iff m\le na+1<nb,\\
    \implies na<m<nb.\qedhere
  \end{gather*}
\end{proof}
